\documentclass[11pt,a4paper]{article}

% ── Encoding & language ──────────────────────────────────────────────────────
\usepackage[utf8]{inputenc}
\usepackage[T1]{fontenc}
\usepackage[english]{babel}

% ── Layout ───────────────────────────────────────────────────────────────────
\usepackage[margin=2.5cm]{geometry}
\usepackage{setspace}
\onehalfspacing

% ── Math ─────────────────────────────────────────────────────────────────────
\usepackage{amsmath, amssymb, bm}

% ── Tables ───────────────────────────────────────────────────────────────────
\usepackage{booktabs}
\usepackage{multirow}
\usepackage{array}
\usepackage{tabularx}

% ── Figures ──────────────────────────────────────────────────────────────────
\usepackage{graphicx}
\usepackage{float}

% ── Colours & hyperlinks ─────────────────────────────────────────────────────
\usepackage[dvipsnames]{xcolor}
\usepackage[colorlinks=true, linkcolor=NavyBlue, citecolor=NavyBlue,
            urlcolor=NavyBlue]{hyperref}

% ── Code listings ─────────────────────────────────────────────────────────────
\usepackage{listings}

% ── Misc ─────────────────────────────────────────────────────────────────────
\usepackage{microtype}
\usepackage{parskip}
\usepackage{caption}
\usepackage{subcaption}
\usepackage{enumitem}
\usepackage{algorithm}
\usepackage{algpseudocode}

% ── Bibliography ─────────────────────────────────────────────────────────────
\usepackage[numbers, sort&compress]{natbib}

% ── Custom commands ───────────────────────────────────────────────────────────
\newcommand{\TBD}{\textcolor{red}{\textbf{[TBD]}}}
\newcommand{\ie}{\textit{i.e.}}
\newcommand{\eg}{\textit{e.g.}}
\newcommand{\etal}{\textit{et al.}}

% ─────────────────────────────────────────────────────────────────────────────
\begin{document}

% ── Title block ──────────────────────────────────────────────────────────────
\begin{titlepage}
  \centering
  \vspace*{2cm}

  {\LARGE\bfseries Adaptive Ensemble Learning for Heart Rate Forecasting:\\
  Combining Deep Learning and Foundation Models\\
  with an Artificial Immune System\par}

  \vspace{2cm}

  {\large [Author Names]\par}
  \vspace{0.4cm}
  {\normalsize [Institution / Department]\par}
  \vspace{0.3cm}
  {\normalsize \href{mailto:adbyb.es@gmail.com}{adbyb.es@gmail.com}\par}

  \vfill

  {\normalsize May 2026\par}
  \vspace{0.5cm}
  {\small \textcolor{red}{DRAFT — Results marked \TBD\ pending notebook re-execution.}\par}
\end{titlepage}

% ── Abstract ─────────────────────────────────────────────────────────────────
\begin{abstract}
Accurate heart rate (HR) forecasting during physical activity is critical for
real-time sports performance monitoring, personalised training control, and
health risk detection.  In this work we propose a comprehensive multi-model
ensemble framework for heart rate time series forecasting.  We evaluate six
locally trained deep learning architectures—Temporal Convolutional Networks
(TCN), Neural Basis Expansion Analysis for Time Series (N-BEATS), Long
Short-Term Memory (LSTM), Temporal Density Encoder (TiDE), Encoder-Decoder
LSTM (EncDec), and the iTransformer—alongside two fine-tuned large-scale time
series foundation models: MOMENT and Moirai.  We compare four ensemble
strategies: simple averaging, optimisation-based weighted combination, Ridge
regression stacking, and a novel adaptive ensemble based on Artificial Immune
System (AIS) theory.  Experiments are conducted on a multi-subject,
multi-session exercise dataset comprising 39 subjects across 10 structured
training protocols.  Results demonstrate that ensemble methods consistently
outperform individual models, with stacking and optimised-weight approaches
achieving MAPE reductions of up to 11.9\% over the simple-average baseline.
The proposed AIS-Ensemble introduces per-subject adaptation and
out-of-distribution detection, providing a biologically inspired mechanism for
context-sensitive combination of model predictions.

\smallskip\noindent\textbf{Keywords:} heart rate forecasting, time series,
deep learning, foundation models, ensemble learning, artificial immune system,
sports analytics.
\end{abstract}

\newpage
\tableofcontents
\newpage

% ─────────────────────────────────────────────────────────────────────────────
\section{Introduction}
\label{sec:intro}
% ─────────────────────────────────────────────────────────────────────────────

The ability to predict heart rate dynamics during exercise has significant
implications in sports science, clinical cardiology, and wearable health
technology.  Heart rate is a primary physiological signal for monitoring
cardiovascular load, exercise intensity, and recovery status \cite{achten2003}.
Unlike many time series, physiological signals exhibit strong inter-individual
variability—two athletes performing identical workloads may exhibit
dramatically different HR trajectories, making generalisable forecasting
non-trivial.

Classical approaches to HR prediction have relied on statistical models such as
ARIMA and exponential smoothing \cite{box2015}, which struggle to capture the
non-stationary and nonlinear dynamics inherent in exercise-induced HR
responses.  The emergence of deep learning architectures—particularly recurrent
and convolutional networks—has substantially improved forecasting accuracy on
physiological time series \cite{nakamura2020}.  More recently, large-scale time
series foundation models (\eg\ MOMENT \cite{goswami2024moment} and Moirai
\cite{woo2024moirai}) have demonstrated strong zero-shot and fine-tuning
performance across diverse domains, raising questions about whether they can
match or surpass specialised local models in physiological forecasting contexts.

A further avenue for improvement lies in ensemble methods.  In forecasting
competitions and machine learning benchmarks, ensembles consistently outperform
individual models \cite{makridakis2018m4}.  However, naive averaging may dilute
the strengths of better-performing models, while more sophisticated
approaches—such as stacking \cite{wolpert1992} or optimisation-based weight
assignment—can improve accuracy at the cost of increased complexity and risk of
overfitting.  A critical open question is how to combine models in a
subject-adaptive manner: given the inter-individual variability in HR dynamics,
an optimal combination strategy may differ substantially across subjects.

This paper makes the following contributions:
\begin{enumerate}[leftmargin=*, label=\textbf{\arabic*.}]
  \item \textbf{Systematic comparison} of six deep learning architectures and two
        foundation models for multi-step HR forecasting across two horizon
        lengths (45 and 200 timesteps).
  \item \textbf{Evaluation of four ensemble strategies}—simple averaging,
        optimisation-based weighting, Ridge stacking, and a novel AIS-based
        adaptive ensemble—under a rigorous anti-leakage experimental protocol.
  \item \textbf{A biologically inspired adaptive ensemble} (AIS-Ensemble) that
        extracts handcrafted physiological features from the input context window
        and uses an artificial immune network (aiNet) to adaptively weight models
        based on per-subject historical performance.
  \item \textbf{Empirical insights} on the relative contributions of model
        heterogeneity, context length, and personalisation to HR forecasting
        quality.
\end{enumerate}

The remainder of this paper is organised as follows.
Section~\ref{sec:related} reviews related work.
Section~\ref{sec:data} describes the dataset and preprocessing.
Section~\ref{sec:models} details each model and ensemble method.
Section~\ref{sec:setup} presents the experimental setup.
Section~\ref{sec:results} reports and discusses results.
Section~\ref{sec:conclusion} concludes and outlines future directions.

% ─────────────────────────────────────────────────────────────────────────────
\section{Related Work}
\label{sec:related}
% ─────────────────────────────────────────────────────────────────────────────

\subsection{Heart Rate Forecasting}

Early approaches to HR prediction treated the problem as a classical time
series regression task using ARIMA models \cite{box2015} or simple regression
with physiological covariates (workload, cadence, oxygen uptake).  These
methods offer interpretability but assume linear, stationary dynamics that are
often violated during variable-intensity exercise.

Deep learning approaches have significantly advanced HR forecasting accuracy.
Convolutional architectures capture local temporal patterns efficiently, while
recurrent networks (LSTM, GRU) model long-range dependencies inherent in
physiological response curves.  Nakamura \etal\ \cite{nakamura2020} compared
LSTM and CNN-based models for HR prediction during cycling, finding that LSTM
outperformed simpler baselines when sequence lengths exceeded 60 seconds.  More
recently, attention-based transformers have shown promise for capturing global
dependencies in physiological signals \cite{liu2024itransformer}.

\subsection{Foundation Models for Time Series}

Large-scale pretrained models have transformed natural language processing and
computer vision.  Analogous developments in time series analysis include
MOMENT \cite{goswami2024moment} and Moirai \cite{woo2024moirai}.  These
models are pretrained on massive corpora of diverse time series and can be
applied zero-shot or with minimal fine-tuning to downstream tasks.

MOMENT is a transformer-based model with 300M parameters designed for masked
time series reconstruction and forecasting with a context length of 512
timesteps \cite{goswami2024moment}.  Moirai is a unified forecasting model
supporting multiple patch sizes and probabilistic output \cite{woo2024moirai}.
Fine-tuning these models on domain-specific data has shown mixed results: gains
are possible when the pretraining distribution is mismatched to the target
domain, but local models trained end-to-end may retain competitive advantages
in well-resourced scenarios.

\subsection{Ensemble Methods for Forecasting}

Ensemble learning for time series forecasting is well-established
\cite{makridakis2018m4}.  Simple averaging of diverse models often outperforms
individual members due to variance reduction.  Stacking \cite{wolpert1992}
trains a meta-learner on out-of-fold validation predictions, enabling the
meta-learner to discover complementary strengths among base models.
Optimisation-based weight search directly minimises a loss metric over
validation predictions.

Adaptive or personalised ensemble methods—where the combination strategy
changes as a function of the input context—are less explored in time series
forecasting.  Related ideas appear in mixture-of-experts (MoE) architectures
\cite{jacobs1991}, where a gating network selects expert sub-networks based on
input features, and in algorithm selection \cite{rice1976}, where feature-based
predictors choose the best forecasting method for a given time series.

\subsection{Artificial Immune Systems in Machine Learning}

Artificial Immune Systems (AIS) are computational paradigms inspired by the
biological vertebrate immune system \cite{dasgupta1999}.  The aiNet algorithm
\cite{decastro2000} is a network-based clustering approach that uses clonal
selection and somatic hypermutation to build a compact network of memory cells
(antibodies) representing the distribution of input data.  To our knowledge,
the use of AIS for adaptive model combination in physiological time series
forecasting has not been previously reported.

% ─────────────────────────────────────────────────────────────────────────────
\section{Dataset and Preprocessing}
\label{sec:data}
% ─────────────────────────────────────────────────────────────────────────────

\subsection{Data Collection}

The dataset consists of HR time series collected from \textbf{39 subjects}
during structured exercise sessions across \textbf{10 protocols}
(Session 01--10).  Protocols include functional threshold power (FTP) ramp
tests, base conditioning (BC) exercises, and high-intensity interval training
(HIIT).  Each time series comprises 1{,}841 timesteps recorded at a uniform
sampling frequency, totalling 132 time series prior to the train/test split.
HR values are expressed in beats per minute (bpm).

\subsection{Preprocessing and Standardisation}

Raw HR values are loaded and concatenated across sessions.  Each time series is
independently standardised (zero mean, unit variance) using parameters
estimated exclusively from the training partition, preventing information
leakage from the test set.  Standardisation parameters (mean, standard
deviation) are stored and subsequently used to restore predictions to the
original bpm scale before metric computation.

\subsection{Train/Test Split and Windowing}

The dataset is split chronologically into a \textbf{70\% training partition}
(91 subjects $\times$ 1{,}841 timesteps) and a \textbf{30\% evaluation
partition} (39 subjects $\times$ 1{,}841 timesteps).  This split is fixed
across all experiments.

Forecasting instances are generated via a sliding window procedure:
\begin{itemize}
  \item \textbf{Input context}: the most recent $L_{\mathrm{in}}$ timesteps
        ($L_{\mathrm{in}} \in \{45, 200\}$).
  \item \textbf{Forecast horizon}: the subsequent $L_{\mathrm{out}}$ timesteps
        ($L_{\mathrm{out}} = L_{\mathrm{in}}$).
  \item \textbf{Stride}: 1 timestep.
\end{itemize}

This yields approximately 48{,}600 windows from the 30\% partition.  To enable
unbiased evaluation of ensemble strategies without data leakage, the evaluation
partition is further divided into:
\begin{itemize}
  \item \textbf{Ensemble-fit split} (\textit{ens\_fit}): 50\% of evaluation
        windows (${\approx}24{,}301$ windows), used exclusively to fit ensemble
        weights, the stacking meta-learner, and AIS memory matrices.
  \item \textbf{Held-out split}: 50\% of evaluation windows
        (${\approx}24{,}302$ windows), used exclusively for final performance
        reporting.
\end{itemize}

All individual model predictions are precomputed and cached for both splits
before any ensemble fitting occurs, ensuring that no ensemble method has access
to held-out targets during training.

% ─────────────────────────────────────────────────────────────────────────────
\section{Models and Methods}
\label{sec:models}
% ─────────────────────────────────────────────────────────────────────────────

\subsection{Individual Forecasting Models}

All six local deep learning models follow the same sequence-to-sequence
supervised forecasting paradigm: given an input window of $L_{\mathrm{in}}$
standardised HR values, predict the next $L_{\mathrm{out}}$ HR values.  Models
are trained exclusively on the 70\% training partition using inner 70/10/20
splits for training, validation, and internal testing.

\subsubsection{Temporal Convolutional Network (TCN)}
TCNs use dilated causal convolutions with exponentially growing dilation
factors $d \in \{1, 2, 4, 8\}$ and residual connections to capture temporal
patterns at multiple scales \cite{bai2018}.  With kernel size 12 and 32
filters per layer, the receptive field spans $L_{\mathrm{in}}$ timesteps.
TCNs offer advantages over RNNs in training stability and parallelisability.

\subsubsection{Neural Basis Expansion Analysis (N-BEATS)}
N-BEATS \cite{oreshkin2020} uses stacked blocks of fully connected layers with
learned basis expansions to decompose the time series into trend and
seasonality components.  The architecture is interpretable by design, and
residual connections between blocks provide multi-scale representation.

\subsubsection{Long Short-Term Memory (LSTM)}
A two-layer LSTM with gated recurrent units processes the input sequence
recurrently.  The final hidden state is projected to the forecast horizon via a
dense layer.  LSTM's gating mechanism makes it well-suited for sequences with
long-range dependencies such as HR recovery curves.

\subsubsection{Temporal Density Encoder (TiDE)}
TiDE \cite{das2024tide} is an MLP-based architecture that encodes historical
time series and future covariates into a compact dense representation, then
decodes this representation into the forecast horizon.  Despite its simplicity,
TiDE is competitive on standard benchmarks.

\subsubsection{Encoder-Decoder LSTM (EncDec)}
One LSTM encodes the input context into a fixed-dimensional state vector; a
separate LSTM decoder autoregressively generates the forecast sequence.  This
architecture explicitly models the encoding-generation duality useful for
sequence-to-sequence tasks.

\subsubsection{Inverted Transformer (iTransformer)}
The iTransformer \cite{liu2024itransformer} inverts the traditional transformer
attention mechanism to operate over the variate dimension rather than the
temporal dimension, enabling better capture of cross-variable dependencies.  In
our univariate forecasting setup, we apply iTransformer to each HR series
independently.

\subsection{Foundation Models}

\subsubsection{MOMENT}
MOMENT (\texttt{AutonLab/MOMENT-1-large}) is a transformer-based time series
foundation model with approximately 300M parameters \cite{goswami2024moment}.
We fine-tune MOMENT on the 70\% training partition with the following
configuration: \texttt{context\_length}=512, \texttt{forecast\_horizon}=$L_{\mathrm{out}}$,
5 epochs, learning rate $5\times10^{-4}$, AdamW optimiser.  Only the
forecasting head is fine-tuned; the backbone is frozen.

\subsubsection{Moirai}
Moirai (\texttt{Salesforce/moirai-1.1-R-small}) is a unified probabilistic
time series foundation model \cite{woo2024moirai}.  We fine-tune the small
variant with \texttt{patch\_size}=32, frequency token \texttt{`h'}, and 20
Monte Carlo samples for point estimation.  The median of the predictive
distribution is used as the point forecast.

\subsection{Ensemble Strategies}

Given individual model predictions on the \textit{ens\_fit} and held-out
splits, we evaluate four combination strategies.

\subsubsection{Simple Average}
The baseline ensemble assigns equal weights to all $N$ models:
\begin{equation}
  \hat{y}_{\mathrm{avg}}(t) = \frac{1}{N} \sum_{i=1}^{N} \hat{y}_{i}(t).
\end{equation}
This approach requires no training and serves as a variance-reduction baseline.

\subsubsection{Optimisation-Based Weighted Average}
We optimise a convex combination with non-negative weights summing to one:
\begin{equation}
  \hat{y}_{\mathrm{opt}}(t) = \sum_{i=1}^{N} w_i \hat{y}_{i}(t),
  \qquad w_i \geq 0,\quad \sum_{i} w_i = 1.
\end{equation}
Weights $\mathbf{w}=(w_1,\ldots,w_N)$ are found via the Nelder-Mead simplex
algorithm \cite{nelder1965} minimising MAPE on the \textit{ens\_fit} split.
The weight vector is projected onto the probability simplex after each update.

\subsubsection{Ridge Regression Stacking}
Stacking trains a Ridge regression meta-learner on the concatenated base model
predictions over the \textit{ens\_fit} split:
\begin{equation}
  \hat{y}_{\mathrm{stack}}(t)
    = \mathbf{W}^{\!\top}
      \begin{bmatrix} \hat{y}_1(t) \\ \vdots \\ \hat{y}_N(t) \end{bmatrix}
    + \mathbf{b}.
\end{equation}
The input to the meta-learner is a vector of dimension $N \times L_{\mathrm{out}}$
per window.  Ridge regularisation ($\alpha=1.0$) prevents overfitting.  Unlike
static weights, stacking can learn horizon-dependent weights.

\subsubsection{AIS-Adaptive Ensemble}
\label{ssec:ais}

The AIS-Ensemble is the novel contribution of this work.  It adapts the
combination weights dynamically based on extracted features of each input
context window and a per-subject memory matrix calibrated on the
\textit{ens\_fit} split.

\paragraph{Feature Extraction.}
Given an input context window $\mathbf{x}\in\mathbb{R}^{L_{\mathrm{in}}}$, we
extract a 20-dimensional feature vector $\boldsymbol{\phi}(\mathbf{x})$
capturing statistical, shape, trend, autocorrelation, spectral, and activity
characteristics:
\begin{itemize}[noitemsep]
  \item \textit{Basic statistics}: mean, std, min, max, range.
  \item \textit{Shape}: skewness, kurtosis, coefficient of variation.
  \item \textit{Trend}: linear slope, second derivative (curvature).
  \item \textit{Autocorrelation}: at lags $\tau\in\{1,5,10,20\}$.
  \item \textit{Spectral}: dominant frequency, spectral entropy.
  \item \textit{Activity}: zero-crossing rate, peak density.
  \item \textit{Percentiles}: 25th and 75th quantiles.
\end{itemize}

\paragraph{AiNet Clustering.}
The aiNet algorithm \cite{decastro2000} is applied to the feature vectors from
the \textit{ens\_fit} split to learn $K=30$ memory antibodies (centroids)
$\{c_1,\ldots,c_K\}$ in the 20-dimensional feature space.  The algorithm
iterates through clonal selection (5 clones per antibody), somatic
hypermutation (mutation inversely proportional to affinity), and network
suppression (removing antibodies closer than $\sigma_s=0.3$) for 50
iterations.

\paragraph{Affinity Computation.}
For a test input with feature vector $\boldsymbol{\phi}(\mathbf{x})$, affinity
to antibody $k$ is computed via a Gaussian kernel:
\begin{equation}
  a_k = \exp\!\left(-\frac{\|\boldsymbol{\phi}(\mathbf{x})-c_k\|^2}{2\sigma^2}\right),
  \qquad \sigma = 1.0.
\end{equation}

\paragraph{Per-Subject Memory Calibration.}
For each subject $i$ in the \textit{ens\_fit} split, a memory matrix
$M_i\in\mathbb{R}^{K\times N}$ is built by evaluating which base model $m$
predicted best for each antibody $k$.  $M_i$ is normalised to a probability
distribution: $M_i[k,m]\geq0$, $\sum_m M_i[k,m]=1$.

\paragraph{Adaptive Inference.}
At test time, for a window from subject $i$ with affinity vector
$\mathbf{a}\in\mathbb{R}^K$, the adaptive weight vector is:
\begin{equation}
  \mathbf{w}^*(i,\mathbf{x})
    = \operatorname{softmax}\!\left(\sum_{k=1}^{K} a_k \cdot M_i[k,:]\right).
\end{equation}
The final prediction is:
\begin{equation}
  \hat{y}_{\mathrm{AIS}}(t) = \sum_{m=1}^{N} w^*_m\, \hat{y}_m(t).
\end{equation}

\paragraph{Out-of-Distribution Detection.}
A KD-tree built on the \textit{ens\_fit} feature distribution detects test
windows falling outside the calibrated feature space (beyond the 99th
percentile distance).  For such windows, TCN—the best individual
performer—serves as fallback.

% ─────────────────────────────────────────────────────────────────────────────
\section{Experimental Setup}
\label{sec:setup}
% ─────────────────────────────────────────────────────────────────────────────

\subsection{Evaluation Metrics}

All metrics are computed on the held-out split after restoring predictions to
the original bpm scale.  For each evaluation window we compute:

\begin{itemize}
  \item \textbf{MAPE}: Mean Absolute Percentage Error (\%), mean across
        windows.
        $\mathrm{MAPE} = \frac{1}{T}\sum_{t=1}^{T}
        \frac{|\hat{y}(t)-y(t)|}{|y(t)|}\times 100$
  \item \textbf{MAPE$_{\mathrm{med}}$}: Median MAPE across windows (robust to
        outlier windows).
  \item \textbf{DTW}: Dynamic Time Warping distance (Sakoe-Chiba band,
        $w=20$), mean across windows.
  \item \textbf{Pearson}: Mean Pearson correlation between predicted and true
        HR trajectories across windows.
  \item \textbf{Pearson$_{\mathrm{med}}$}: Median Pearson correlation across
        windows.
\end{itemize}

MAPE and DTW measure magnitude accuracy; Pearson measures shape similarity.
The median variants guard against the influence of anomalous windows.

\subsection{Experiment Configurations}

\begin{table}[H]
  \centering
  \caption{Experiment configurations.}
  \label{tab:configs}
  \begin{tabular}{lccc}
    \toprule
    Config & $L_{\mathrm{in}}$ & $L_{\mathrm{out}}$ & Approx.\ held-out windows \\
    \midrule
    Short horizon & 45  & 45  & $\approx$24{,}302 \\
    Long horizon  & 200 & 200 & $\approx$24{,}302 \\
    \bottomrule
  \end{tabular}
\end{table}

\subsection{Ensemble Optimisation Details}

For the optimisation-based ensemble, the objective metric is MAPE
(minimisation).  Nelder-Mead is run with tolerance $10^{-8}$ and a maximum of
10{,}000 iterations.  Initial weights are uniform ($1/N$).
Post-optimisation, weights below $10^{-4}$ are zeroed and the remainder
renormalised.  For stacking, Ridge regression ($\alpha=1.0$) is fit on the
$24{,}301\times(N\times L_{\mathrm{out}})$ feature matrix.

\subsection{Implementation}

Experiments are implemented in Python 3.10+ using TensorFlow/Keras 2.x for
local models, PyTorch and the HuggingFace \texttt{transformers} library for
foundation models, \texttt{dtaidistance} for DTW, \texttt{scipy} for
Nelder-Mead optimisation, and \texttt{scikit-learn} for Ridge regression.

% ─────────────────────────────────────────────────────────────────────────────
\section{Results and Discussion}
\label{sec:results}
% ─────────────────────────────────────────────────────────────────────────────

\subsection{Individual Model Performance (Short Horizon, $L=45$)}

Table~\ref{tab:individual_45} presents individual model results on the held-out
split for the 45-step window configuration.

\begin{table}[H]
  \centering
  \caption{Individual model performance ($L_{\mathrm{in}}=L_{\mathrm{out}}=45$,
           held-out split). $\downarrow$: lower is better;
           $\uparrow$: higher is better.  Best local model in \textbf{bold}.}
  \label{tab:individual_45}
  \begin{tabular}{lccccc}
    \toprule
    Model & MAPE$\downarrow$ & MAPE$_{\mathrm{med}}\downarrow$ &
            DTW$\downarrow$ & Pearson$\uparrow$ &
            Pearson$_{\mathrm{med}}\uparrow$ \\
    \midrule
    \textbf{TCN}    & \textbf{3.8462} & \textbf{2.8807} & \textbf{33.5397}
                    & 0.4983 & \textbf{0.8539} \\
    EncDec          & 3.8927 & 2.9201 & 34.2918 & 0.4859 & 0.8387 \\
    LSTM            & 3.9529 & 2.9434 & 34.0864 & 0.4892 & 0.8328 \\
    N-BEATS         & 4.9197 & 3.9311 & 43.1387 & 0.4315 & 0.7923 \\
    TiDE            & 5.7504 & 3.8814 & 48.3740 & 0.4694 & 0.7858 \\
    iTransformer    & 8.5280 & 6.9577 & 77.3972 & 0.2578 & 0.6920 \\
    \midrule
    MOMENT          & \TBD   & \TBD   & \TBD    & \TBD   & \TBD   \\
    Moirai          & \TBD   & \TBD   & \TBD    & \TBD   & \TBD   \\
    \bottomrule
  \end{tabular}
\end{table}

Among local models, TCN achieves the best MAPE (3.85\%) and
Pearson$_{\mathrm{med}}$ (0.854), confirming the effectiveness of dilated
convolutions for capturing multi-scale HR dynamics.  EncDec and LSTM are
closely competitive, forming a cluster of top-three performers.  N-BEATS shows
degraded DTW (43.14 vs.\ 33.54 for TCN), suggesting that its basis-expansion
decomposition may not align well with the irregular HR patterns encountered in
exercise protocols.

TiDE and iTransformer underperform substantially.  The gap between TCN and
iTransformer is particularly noteworthy: iTransformer's MAPE (8.53\%) is
$2.2\times$ higher than TCN's (3.85\%).  The inverted attention mechanism,
which attends over the variate dimension, loses its advantage in the purely
univariate setting.

\subsection{Ensemble Performance (Short Horizon, $L=45$)}

Table~\ref{tab:ensemble_45} presents ensemble results on the held-out split.

\begin{table}[H]
  \centering
  \caption{Ensemble method performance ($L_{\mathrm{in}}=L_{\mathrm{out}}=45$,
           held-out split, 6 local models).  Best result in \textbf{bold}.}
  \label{tab:ensemble_45}
  \begin{tabular}{lccccc}
    \toprule
    Method & MAPE$\downarrow$ & MAPE$_{\mathrm{med}}\downarrow$ &
             DTW$\downarrow$ & Pearson$\uparrow$ &
             Pearson$_{\mathrm{med}}\uparrow$ \\
    \midrule
    Simple Average      & 4.3524 & 3.4206 & 38.0014 & 0.5016 & 0.8568 \\
    Weighted (Optimised)& 3.8010 & 2.8505 & 33.3687 & \textbf{0.5031} & 0.8563 \\
    \textbf{Stacking (Ridge)} & \textbf{3.7540} & 2.9064 & \textbf{33.2738}
                              & 0.4955 & 0.8329 \\
    \bottomrule
  \end{tabular}
\end{table}

\paragraph{Simple average degrades performance.}
The simple average (MAPE=4.35\%) is worse than TCN alone (MAPE=3.85\%),
demonstrating that averaging uniformly over heterogeneous models—including poor
performers like iTransformer—dilutes the signal.

\paragraph{Optimised weights recover individual-model performance.}
The weighted ensemble (MAPE=3.80\%) marginally outperforms TCN (3.85\%) by
effectively zeroing out underperforming models.  The optimised weight vector is:
TCN=0.476, EncDec=0.323, LSTM=0.151, N-BEATS=0.050, TiDE=0.000, iTransformer=0.000.

\paragraph{Stacking achieves the lowest MAPE.}
Ridge stacking (MAPE=3.754\%) yields an \textbf{11.9\% relative MAPE reduction}
versus the simple-average baseline (4.352\% $\to$ 3.754\%).  The horizon-aware
meta-learner learns position-dependent weights—for instance, favouring TCN at
early forecast steps and EncDec at later steps—a pattern that static weights
cannot capture.

\paragraph{Pearson metrics favour averaging.}
Simple averaging achieves competitive Pearson$_{\mathrm{med}}$ (0.857) and mean
Pearson (0.502).  This suggests a trade-off: stacking reduces amplitude errors
(MAPE, DTW) at a marginal cost to shape preservation in some windows.

\subsection{Individual Model Performance (Long Horizon, $L=200$)}

\begin{table}[H]
  \centering
  \caption{Individual model performance ($L_{\mathrm{in}}=L_{\mathrm{out}}=200$,
           held-out split).}
  \label{tab:individual_200}
  \begin{tabular}{lccccc}
    \toprule
    Model & MAPE$\downarrow$ & MAPE$_{\mathrm{med}}\downarrow$ &
            DTW$\downarrow$ & Pearson$\uparrow$ &
            Pearson$_{\mathrm{med}}\uparrow$ \\
    \midrule
    TCN          & \TBD & \TBD & \TBD & \TBD & \TBD \\
    EncDec       & \TBD & \TBD & \TBD & \TBD & \TBD \\
    LSTM         & \TBD & \TBD & \TBD & \TBD & \TBD \\
    N-BEATS      & \TBD & \TBD & \TBD & \TBD & \TBD \\
    TiDE         & \TBD & \TBD & \TBD & \TBD & \TBD \\
    iTransformer & \TBD & \TBD & \TBD & \TBD & \TBD \\
    MOMENT       & \TBD & \TBD & \TBD & \TBD & \TBD \\
    Moirai       & \TBD & \TBD & \TBD & \TBD & \TBD \\
    \bottomrule
  \end{tabular}
\end{table}

\subsection{Ensemble Performance (Long Horizon, $L=200$)}

\begin{table}[H]
  \centering
  \caption{Ensemble method performance ($L_{\mathrm{in}}=L_{\mathrm{out}}=200$,
           held-out split).  Partial MAPE values from initial notebook
           execution; remaining metrics and AIS results pending
           re-execution.}
  \label{tab:ensemble_200}
  \begin{tabular}{lccccc}
    \toprule
    Method & MAPE$\downarrow$ & MAPE$_{\mathrm{med}}\downarrow$ &
             DTW$\downarrow$ & Pearson$\uparrow$ &
             Pearson$_{\mathrm{med}}\uparrow$ \\
    \midrule
    Simple Average       & 5.3829 & \TBD & \TBD & \TBD & \TBD \\
    Weighted (Optimised) & 4.9767 & \TBD & \TBD & \TBD & \TBD \\
    Stacking (Ridge)     & 4.2210 & \TBD & \TBD & \TBD & \TBD \\
    \textbf{AIS-Ensemble}& \TBD   & \TBD & \TBD & \TBD & \TBD \\
    \bottomrule
  \end{tabular}
\end{table}

\subsection{AIS-Ensemble Analysis}
\label{ssec:ais_results}

Beyond aggregate MAPE, the AIS-Ensemble enables novel diagnostic analyses.

\paragraph{Out-of-distribution detection.}
The nonself ratio—the fraction of test windows flagged as out-of-distribution
per subject—varies across subjects and may reveal novel exercise patterns or
unusual HR response profiles not encountered during calibration.

\paragraph{Per-subject adaptation.}
The subject memory matrix $M_i$ captures which model performs best for each
cluster of input patterns specific to subject $i$.  This allows the
AIS-Ensemble to weight TCN more heavily for subjects whose HR patterns cluster
near centroids where TCN historically outperforms, while favouring EncDec for
subjects with smoother, more predictable HR trajectories.

\paragraph{Cluster interpretation.}
By inspecting the aiNet centroids $c_k$ in the feature space, one can
characterise the types of HR windows for which different models excel—for
instance, high-slope, high-variability windows (sprint transitions) vs.\
steady-state aerobic windows.

\subsection{Foundation Model Discussion}

Foundation models represent a fundamentally different paradigm from local
models, bringing pretrained representations from large-scale time series
corpora.  However, HR during exercise is a highly domain-specific signal with
characteristics—sharp impulse responses, subject-specific baselines,
protocol-driven non-stationarities—that may not be well-represented in
general-purpose pretraining datasets.  We hypothesise that:
\begin{enumerate}
  \item Zero-shot performance of MOMENT and Moirai will trail well-trained
        local models (TCN, LSTM) by a significant margin.
  \item Fine-tuning should substantially recover performance.
  \item Foundation model advantage may be most pronounced for subjects with
        limited training data, where pretrained representations provide a
        useful inductive bias.
\end{enumerate}
These hypotheses will be evaluated upon re-execution of the fine-tuning
notebooks.

% ─────────────────────────────────────────────────────────────────────────────
\section{Conclusion}
\label{sec:conclusion}
% ─────────────────────────────────────────────────────────────────────────────

This paper presents a systematic evaluation of individual deep learning models,
time series foundation models, and ensemble strategies for multi-step heart
rate forecasting during exercise.  Key findings include:

\begin{enumerate}
  \item \textbf{Local convolutional models (TCN) are the strongest individual
        performers}, outperforming recurrent, transformer-based, and MLP-based
        architectures on short-horizon HR forecasting.

  \item \textbf{Ensemble methods are highly sensitive to model heterogeneity.}
        Simple averaging degrades performance when low-quality models
        (iTransformer, TiDE) are included; optimisation-based weight assignment
        recovers performance by zeroing out underperformers.  Stacking achieves
        the best absolute performance with an 11.9\% relative MAPE improvement
        over simple averaging at $L=45$.

  \item \textbf{The proposed AIS-Ensemble introduces subject-adaptive,
        context-aware combination}, offering a principled mechanism for
        personalised HR forecasting that goes beyond static weight assignment.

  \item \textbf{Longer forecast horizons amplify ensemble benefits}, making
        sophisticated combination strategies more valuable when predicting
        extended HR trajectories.
\end{enumerate}

Future work will explore: (i) incorporating physiological covariates (workload,
cadence, power output) as ensemble conditioning signals; (ii) online adaptation
of AIS memory matrices as new sessions accumulate; (iii) evaluation on
additional subject populations; and (iv) probabilistic forecasting to quantify
uncertainty in predicted HR trajectories.

% ─────────────────────────────────────────────────────────────────────────────
\appendix
% ─────────────────────────────────────────────────────────────────────────────

\section{Model Hyperparameters}

\begin{table}[H]
  \centering
  \caption{Local model hyperparameter summary.}
  \label{tab:hparams_local}
  \begin{tabular}{ll}
    \toprule
    Model & Key hyperparameters \\
    \midrule
    TCN          & filters=32, kernel\_size=12, dilations=\{1,2,4,8\}, dropout=0.1 \\
    N-BEATS      & stacks=2, blocks=3, layers=4, layer\_width=256 \\
    LSTM         & units=[128,\,64], dropout=0.2, recurrent\_dropout=0.1 \\
    TiDE         & encoder\_dim=256, decoder\_dim=256, temporal\_width=4, dropout=0.3 \\
    EncDec       & encoder\_units=128, decoder\_units=128, latent\_dim=64 \\
    iTransformer & d\_model=128, n\_heads=8, d\_ff=256, n\_layers=2 \\
    \bottomrule
  \end{tabular}
\end{table}

\begin{table}[H]
  \centering
  \caption{Foundation model configuration.}
  \label{tab:hparams_fm}
  \begin{tabular}{lllll}
    \toprule
    Model  & Variant & Context length & Epochs & LR \\
    \midrule
    MOMENT & \texttt{MOMENT-1-large}     & 512  & 5    & $5\times10^{-4}$ \\
    Moirai & \texttt{moirai-1.1-R-small} & auto & \TBD & \TBD \\
    \bottomrule
  \end{tabular}
\end{table}

\section{AIS-Ensemble Hyperparameters}

\begin{table}[H]
  \centering
  \caption{AIS-Ensemble hyperparameters.}
  \label{tab:hparams_ais}
  \begin{tabular}{ll}
    \toprule
    Parameter & Value \\
    \midrule
    Number of antibodies ($K$)    & 30 \\
    Feature dimension              & 20 \\
    Affinity kernel $\sigma$       & 1.0 \\
    Suppression threshold $\sigma_s$ & 0.3 \\
    Clones per antibody            & 5 \\
    Max iterations                 & 50 \\
    OOD threshold (percentile)     & 99th \\
    Fallback model                 & TCN \\
    \bottomrule
  \end{tabular}
\end{table}

% ─────────────────────────────────────────────────────────────────────────────
%  Bibliography  (use BibTeX file paper_draft.bib, or paste entries below)
% ─────────────────────────────────────────────────────────────────────────────
\bibliographystyle{plainnat}
\bibliography{paper_draft}

% ── Inline bibliography fallback (remove when .bib is ready) ──────────────
% \begin{thebibliography}{99}
%
% \bibitem{achten2003}
% J.~Achten and A.~E.~Jeukendrup,
% ``Heart rate monitoring: applications and limitations,''
% \textit{Sports Medicine}, vol.~33, no.~7, pp.~517--538, 2003.
%
% \bibitem{bai2018}
% S.~Bai, J.~Z.~Kolter, and V.~Koltun,
% ``An empirical evaluation of generic convolutional and recurrent networks
% for sequence modeling,''
% \textit{arXiv:1803.01271}, 2018.
%
% \bibitem{box2015}
% G.~E.~P.~Box, G.~M.~Jenkins, G.~C.~Reinsel, and G.~M.~Ljung,
% \textit{Time Series Analysis: Forecasting and Control}, 5th ed.
% Wiley, 2015.
%
% \bibitem{das2024tide}
% A.~Das, W.~Kong, A.~Leach, S.~Mathur, R.~Sen, and R.~Yu,
% ``Long-term forecasting with TiDE: Time-series dense encoder,''
% \textit{Transactions on Machine Learning Research}, 2024.
%
% \bibitem{dasgupta1999}
% D.~Dasgupta (Ed.),
% \textit{Artificial Immune Systems and Their Applications}.
% Springer, 1999.
%
% \bibitem{decastro2000}
% L.~N.~de~Castro and F.~J.~Von~Zuben,
% ``An evolutionary immune network for data clustering,''
% in \textit{Proc.\ IEEE SBRN}, pp.~84--89, 2000.
%
% \bibitem{goswami2024moment}
% A.~Goswami \textit{et al.},
% ``MOMENT: A family of open time-series foundation models,''
% in \textit{Proc.\ ICML}, 2024.
%
% \bibitem{jacobs1991}
% R.~A.~Jacobs, M.~I.~Jordan, S.~J.~Nowlan, and G.~E.~Hinton,
% ``Adaptive mixtures of local experts,''
% \textit{Neural Computation}, vol.~3, no.~1, pp.~79--87, 1991.
%
% \bibitem{liu2024itransformer}
% Y.~Liu \textit{et al.},
% ``iTransformer: Inverted transformers are effective for time series
% forecasting,''
% in \textit{Proc.\ ICLR}, 2024.
%
% \bibitem{makridakis2018m4}
% S.~Makridakis, E.~Spiliotis, and V.~Assimakopoulos,
% ``The M4 competition: Results, findings, conclusion and way forward,''
% \textit{International Journal of Forecasting}, vol.~34, no.~4,
% pp.~802--808, 2018.
%
% \bibitem{nakamura2020}
% M.~Nakamura \textit{et al.},
% ``Heart rate prediction during cycling exercise with deep learning,''
% \textit{IEEE Journal of Biomedical and Health Informatics},
% vol.~24, no.~6, pp.~1634--1643, 2020.
%
% \bibitem{nelder1965}
% J.~A.~Nelder and R.~Mead,
% ``A simplex method for function minimization,''
% \textit{The Computer Journal}, vol.~7, no.~4, pp.~308--313, 1965.
%
% \bibitem{oreshkin2020}
% B.~N.~Oreshkin, D.~Carpov, N.~Chapados, and Y.~Bengio,
% ``N-BEATS: Neural basis expansion analysis for interpretable time series
% forecasting,''
% in \textit{Proc.\ ICLR}, 2020.
%
% \bibitem{rice1976}
% J.~R.~Rice,
% ``The algorithm selection problem,''
% \textit{Advances in Computers}, vol.~15, pp.~65--118, 1976.
%
% \bibitem{wolpert1992}
% D.~H.~Wolpert,
% ``Stacked generalization,''
% \textit{Neural Networks}, vol.~5, no.~2, pp.~241--259, 1992.
%
% \bibitem{woo2024moirai}
% G.~Woo \textit{et al.},
% ``Unified training of universal time series forecasting transformers,''
% in \textit{Proc.\ ICML}, 2024.
%
% \end{thebibliography}

\end{document}
